\section{Outstanding analytical tasks}\label{sec:outstanding}

The manuscript supplies exact local identities and identifies useful
literature inputs. The missing step is a quantitative restriction
from the canonical source that the colliding controls do not share.
The following tasks distinguish that step from further changes of
coordinates or from an unproved extension of a statistical theorem.

\subsection{Acquire a coupled coefficient estimate}
The natural local quantities are the critical amplitude \(a\),
the drift \(\beta=3\ell\), and the relative curvature
\(\gamma=6\kappa-3\ell^2\). Their sign structure is now clear:
the conditional real-zero benchmark and an actual colliding control
both have \(\gamma<0\). A useful next estimate must instead
control magnitude, variation, or a coupled relation to the critical
value on a specified time-height region. A collection of separate
upper bounds on jets is insufficient without the appropriate
denominator or boundary margin.

The two exact acquisition interfaces are the signed theta moments
and primitive derivatives in Section~\ref{sec:coefficients}, and
the contour formula for \(U'/U\). Before estimating either,
specify the consumer, the region and the desired gain. The moment
route must pay its coupled cancellation and any curvature lower
bound; the contour route must pay the nonzero boundary margin,
cluster count, source errors and derivatives. The existing
same-source scalar baseline remains the comparison. No advantage
is established just by moving the same unknown quantity to a
primitive or contour.

\subsection{Resolve the statistical gap at sparse exceptions}
An unconditional complex-difference correlation theorem already
exists. Reproving its definition or positivity would not address
the upper-bound problem. Section~\ref{sec:complex-statistics}
shows two separate limitations: a density-one result permits
exceptions, and fixed-positive-time bulk statistics see an
eventually real, regularly spaced tail even when
\(0<\lambda<\Lambda\).

A statistical continuation should therefore select one of two
substantive targets: a localized estimate that detects even a
single exceptional cluster, or a proved amplification from one
exception to an arithmetic statistic with a contradictory bound.
For the first, \eqref{eq:transverse-balance} isolates the exterior
term requiring canonical control; its monotonicity in a polynomial
model alone does not exclude a collision. For the second, the
amplification theorem itself is missing. Neither a new marked
measure nor replacing a complex interaction by its absolute value
supplies that theorem. The full ordinate pair-correlation
conjecture must remain a conjectural input when used.

\subsection{Pay the time-height limit and event coverage}
The normal forms have event-dependent neighborhoods and constants.
A regional result must control remainders, cluster separation and
complete arithmetic and radial tails throughout its declared
region. A global statement needs coverage beyond every finite
height, higher multiplicities and endpoint events, without an
unstated assumption that the threshold is attained at an isolated
finite collision.

The positive-time tail theorem is not uniform as time tends to zero.
Any joint limit \(\lambda=\lambda(T)\downarrow0\) must carry
explicit time-dependent error and exceptional-height bounds.
Such a limit may clarify the transition in bulk statistics, but
it must also retain sensitivity to bounded-height exceptions
if its intended conclusion is \(\Lambda\le0\). An estimate
equivalent to that conclusion should be identified as a
reformulation, rather than counted as an independent bound.

\subsection{Assessment and supporting evidence}
A separate bounded mathematical assessment should check the
coefficient and primitive identities, moving-boundary term,
Laguerre--P\'olya hypotheses, finite-exception bounds, and
transverse-cluster balance, together with the normalizations and
quantifiers of the imported statistical results. The current
analytical and production checks are proof-author checks, not
specialist review. Numerics, if later commissioned, should
evaluate one specified inequality with complete error payments
and remain supporting evidence.

Intrinsic preparation of the theta source and a possible
physical or computational access advantage remain separate
questions. The present operator realization faithfully
transports the supplied source; its exact scalar error comparison
does not decide either question.
